% Exercise 21 (Renewing Commands)

\documentclass[12pt]{scrbook}

% datetime options:
% short: abbreviated format
% nodayofweek: don't display the day of the week
% level: put the "th" at the same level as the text.
% 12hr or 24hr: select 12 or 24 hour format
%\usepackage[short,nodayofweek,level,12hr]{datetime}
%\usepackage[short,nodayofweek,level,24hr]{datetime}
%\usepackage[short,nodayofweek,12hr]{datetime}
%\usepackage[nodayofweek,12hr]{datetime}
\usepackage[T1]{fontenc}
\usepackage[short,12hr]{datetime}
\usepackage{graphicx}
\usepackage[english]{babel}

\pagestyle{headings}

\title{A Simple Document}
\author{Ashwin Menon}

\addto{\captionsenglish}{%
    \renewcommand{\figurename}{Fig}%
    \renewcommand{\tablename}{Tab}%
}

\begin{document}

\maketitle

\frontmatter
\tableofcontents
\listoffigures
\listoftables

\chapter{Summary}
A brief document to
illustrate how to use \LaTeX.

\mainmatter

\chapter{Introduction}
\label{ch:intro}

\section{The First Section}
This is a simple \LaTeX\ document. Here is the first paragraph.
The next chapter is Chapter~\ref{ch:another}
and is on page~\pageref{ch:another}.
The next section is Section~\ref{sec:next}.

\section{The Next Section}
\label{sec:next}

Here is the second paragraph\footnote{with a footnote}. As you
can see it's a rather short paragraph, but not as short as the
previous one. This document was created on: \today\ at \currenttime.

\section{The Shapes Section}
Here are some shapes inline: \includegraphics[height=1ex]{shapes}

Here are some shapes rotated inline: \includegraphics[height=1ex,angle=60]{shapes}

Here are some shapes in a separate paragraph:

\includegraphics{shapes}

Figure~\ref{fig:shapes} shows a picture of shapes in a \emph{figure} environment:

\begin{figure}[htbp]
    \centering
    \includegraphics{shapes}
    \caption{Some shapes}
    \label{fig:shapes}
\end{figure}

\section{The Funky Text Section}
Rotating text: \rotatebox{90}{straight up}
\rotatebox{45}{rising up}
\rotatebox[origin=cB]{45}{rising up} % origin = (c)enter (B)aseline
\rotatebox[units=-360]{45}{falling down}
that's all.

\resizebox{20mm}{1cm}{Squished text}
\resizebox{!}{1cm}{Big text}

\chapter{Another Chapter}
\label{ch:another}

Here's another very interesting chapter.
We're going to put a picture here later.
See Chapter~\ref{ch:intro} for an introduction.

\chapter{Tables}

Table \ref{tab:travel_expenses} shows travel expenses.
\begin{table}[htbp]
    \caption{Travel Expenses}
    \label{tab:travel_expenses}
    \centering
    \begin{tabular}{lrr}
                        & \multicolumn{2}{c}{\bfseries Expenditure (\pounds)} \\
                        & \multicolumn{1}{c}{Year1} & \multicolumn{1}{c}{Year2} \\
    \bfseries Travel    & 100,000 & 110,000 \\
    \bfseries Equipment & 50,000  & 60,000
    \end{tabular}
\end{table}

Table \ref{tab:sample_table} is a sample table.
\begin{table}[htbp]
    \caption{A Sample Table}
    \label{tab:sample_table}
    \centering
    \begin{tabular}{lr}
        Item  & Cost \\
        Video & 8.99 \\
        CD    & 9.99 \\
        DVD   & 15.00 \\
    \end{tabular}
\end{table}

\chapter{Recommended Reading}
For a basic introduction to \LaTeX\ see Lamport~\cite{lamport94}. For more
detailed information about \LaTeX\ and associated applications, consult
Kopka and Daly~\cite{kopka95} or Goossens \emph{et al}~\cite{goossens94}.

\chapter*{Acknowledgements}

I would like to acknowledge all those
very helpful people who have assisted me in my work,

\appendix

\chapter{Tables}

\begin{tabular}{lrr}
                     & \multicolumn{2}{c}{\bfseries Expenditure (\pounds)} \\
                     & \multicolumn{1}{c}{Year1} & \multicolumn{1}{c}{Year2} \\
 \bfseries Travel    & 100,000 & 110,000 \\
 \bfseries Equipment & 50,000  & 60,000
\end{tabular}

\begin{thebibliography}{Goossens 1994}
    \bibitem[Lamport 1994]{lamport94} ''\LaTeX\ : a document preparation system'',
    Leslie Lamport, 2nd edition (updated for \LaTeXe), Addison-Wesley (1994).

    \bibitem[Kopka 1995]{kopka95} ''A Guide to \LaTeX: document preparation for
    beginners and advanced users'', Helmut Kopka and Patrick W. Daly, 
    Addison-Wesley (1995).

    \bibitem[Goossens 1994]{goossens94} ''The \LaTeX\ Companion'', Michel Goossens,
    Frank Mittelbach and Alexander Samarin, Addison-Wesley (1994).
\end{thebibliography}

\end{document}
